\fichatitulo{40 · RISCOS}{Preprint v1 · 2024}{Large Legal Fictions}
{\small Dahl et al. \textit{Large Legal Fictions: Profiling Legal Hallucinations in Large Language Models}. Preprint v1 · 2024. \href{https://arxiv.org/pdf/2401.01301v1}{Fonte consultada}.\par}

\campo{Problema e objetivo}
Com que frequência modelos erram perguntas verificáveis sobre decisões judiciais?

\campo{Principal contribuição · síntese interpretativa}
Avalia alucinações jurídicas e a capacidade de reconhecer limites ou premissas incorretas.

\campo{Método / natureza da evidência}
Perguntas sobre casos federais dos Estados Unidos, incluindo informações específicas e verificáveis e formulações contrafactuais.

\campo{Resultado ou argumento principal}
O resumo da v1 relata taxas de 69\% para GPT-3.5 a 88\% para Llama 2 nas condições avaliadas; não são taxas gerais para todo uso jurídico ou RAG.

\campo{Excertos-chave · seleção literal}
\begin{itemize}
\excerto{1}{specific, verifiable questions}{Resumo.}
\excerto{2}{cannot always predict}{Resumo.}
\end{itemize}

\campo{Limitações e análise crítica}
Análise da ficha: domínio, formulação das perguntas e versões dos modelos delimitam os resultados. Não transferir essas porcentagens para aplicações parlamentares com recuperação.

\campo{Aproveitamento na dissertação · proposta}
Problema e metodologia: incluir verificabilidade, perguntas sem suporte e tratamento de premissas falsas.

\registro{Fonte consultada nesta preparação: Resumo e introdução/tipologia. A bibliografia cita a publicação; ficha usa arXiv v1. Chave: \texttt{dahlEtAl2024largeLegalFictions}.}
